\subsubsection{Gauge symmetries and Abel translations}
\label{sec:gauge-Abel-symmetries}

We interpret the conserved charges and spectral flows as symmetries,
so that their roles in the integral count and the BA phase data can
be distinguished. A conserved function $Q$ generates a canonical
symmetry because $\{Q,H\}=0$ implies that its Hamiltonian flow
commutes with the time evolution. This is the Hamiltonian form of
Noether's correspondence.
For the resonant three-wave system, reduction by mode-phase symmetries
and reconstruction of the geometric and dynamical phases were developed
by Alber, Luther, Marsden and Robbins
\cite[Sections~2--4]{AlberLutherMarsdenRobbins1998}.

For arbitrary $L,M,N$, the power charges in
\cref{eq:canonical-charges} generate the constant reciprocal gauges
\begin{equation}\label{eq:general-mode-gauge-action}
\begin{aligned}
 u_{\ell m}&\longmapsto e^{\theta_\ell^L-\theta_m^M}u_{\ell m},\\
 v_{mn}&\longmapsto e^{\theta_m^M-\theta_n^N}v_{mn},\\
 w_{n\ell}&\longmapsto e^{\theta_n^N-\theta_\ell^L}w_{n\ell},
 \qquad1\leq\ell\leq L,\quad1\leq m\leq M,\quad1\leq n\leq N.
\end{aligned}
\end{equation}
The barred partners have inverse factors. The $\theta$ parameters
are constant complex numbers; their common shift acts trivially,
leaving $d-1$ independent gauges. Each cubic interaction is invariant
because its three exponents sum to zero.

On the regular spectral locus, the ambient integral count separates these gauge generators
from the remaining spectral integrals:
\begin{equation}\label{eq:gauge-translation-count}
 p=(d-1)+G,\qquad d=L+M+N,
 \qquad p=LM+MN+NL,
\end{equation}
where $G$ is the regular spectral genus in
\cref{eq:canonical-genus-preview}. Locally, fixing the $d-1$ independent
charges and quotienting their gauge actions leaves $G$ Hamiltonian
degrees of freedom. Under eigenvector-divisor linearization, the
commuting isospectral flows act as Abel translations
\cite{Beauville1990,Vanhaecke2001}. The autonomous Hamiltonian selects
the time direction within this commuting family; it is included in the
complete integral set of \cref{thm:canonical-independence}.

When phases coincide, the moment matrices
\cref{eq:canonical-block-moment} extend the diagonal gauges to
continuous changes of basis within equal-phase groups.
Their principal-block invariants provide the additional integrals
counted in \cref{eq:canonical-spectral-deficit}.

On a prescribed birational BA realization,
\cref{thm:BA-higher-flows} supplies the translations of its Jacobian
and the accompanying endpoint phases. The power-charge flows change
the multiplicative mode normalizations. These give the two symmetry
actions visible in the BA formulas; the hyperelliptic divisor lift in
\cref{sec:mumford-gauge-lift} makes their distinction explicit.
